\subsection{Hitchin fibers} An important observation made by Hitchin is that the fibers of the Hitchin fibration $f$ can be described in ``abelian'' terms, namely by line bundles on certain finite coverings of $X$. We elaborate on this observation for $G=\GL_{n}$ and $G=\Sp_{2n}$.

\subsubsection{The case of $\GL_{n}$ and the spectral curve}\label{sss:spectral curve} For $a=(a_{1},\cdots, a_{n})\in\calA_{n,\calL}$, one can define a curve $Y_{a}$ equipped with a degree $n$ morphism $\pi_{a}: Y_{a}\to X$. The construction is as follows. The total space of $\calL$ can be written as a relative spectrum over $X$
\begin{equation*}
\Sigma:=\Tot(\calL)=\Spec(\calO\oplus\calL^{\otimes -1}\oplus\calL^{\otimes -2}\oplus\cdots)=\Spec\Sym(\calO\oplus\calL^{\otimes -1})
\end{equation*}
Let $\pi: \Sigma\to X$ be the projection. Consider the map of coherent sheaves on $X$
\begin{equation*}
\iota_{a}: \calL^{\otimes -n}\to \pi_{*}\calO_{\Sigma}=\calO\oplus\calL^{\otimes -1}\oplus\calL^{\otimes -2}\oplus\cdots
\end{equation*}
given in coordinates by $((-1)^{n}a_{n}, (-1)^{n-1}a_{n-1},\cdots, -a_{1},1,0,\cdots)$. By adjunction $\iota_{a}$ induces a map $\iota'_{a}: \pi^{*}\calL^{\otimes -n}\to \calO_{\Sigma}$, whose image we denote by $\calI_{a}$.  Then $\calI_{a}$ is an ideal sheaf on $\Sigma$. Define the {\em spectral curve} $Y_{a}$ to be the closed subscheme of $\Sigma$ defined by $\calI_{a}$
\begin{equation*}
Y_{a}=\Spec(\calO\oplus\calL^{\otimes-1}\oplus\calL^{\otimes-2}\oplus\cdots)/\calI_{a}
\end{equation*}
If we trivialize $\calL$ on some open subset $U\subset X$, and view $a_{i}$ as functions on $U$, then $Y_{a}|_{U}$ is the subscheme of $U\times\AA^{1}$ defined by one equation $y^{n}-a_{1}y^{n-1}+a_{2}y^{n-2}+\cdots +(-1)^{n}a_{n}=0$ (where $y$ is the coordinate on $\AA^{1}$). The projection $\pi_{a}: Y_{a}\to X$ is finite flat of degree $n$. The curve $Y_{a}$ is called the {\em spectral curve} of $a$ since the fibers of $\pi_{a}$ are the roots of the characteristic polynomial $y^{n}-a_{1}y^{n-1}+a_{2}y^{n-2}+\cdots +(-1)^{n}a_{n}$.

When $a\in\Ah_{n,\calL}$,  the curve $Y_{a}$ is reduced and therefore smooth on a Zariski dense open subset, there is a moduli stack $\cPic(Y_{a})$ classifying torsion-free coherent $\calO_{Y_{a}}$-modules that are generically of rank $1$, see \cite{AK}. The usual Picard stack $\Pic(Y_{a})$ classifying line bundles on $Y_{a}$ is an open substack of $\cPic(Y_{a})$, and it acts on $\cPic(Y_{a})$ by tensoring.


\begin{prop}\label{p:cPic}  Suppose $a\in\Ah_{n,\calL}(k)$. Let $\calM_{a}$ be the fiber of $f:\calM_{n,\calL}\to \calA_{n,\calL}$ over $a$. Then there is a canonical isomorphism of stacks
\begin{equation*}
\cPic(Y_{a})\cong\calM_{a}.
\end{equation*}
\end{prop}

\subsubsection{Sketch of proof}\label{sss:pf cPic}
We give the morphism $\cPic(Y_{a})\to \calM_{a}$. For any coherent sheaf $\calF$ on $Y_{a}$, the direct image $\pi_{a,*}\calF$ is a coherent sheaf on $X$ equipped with a map $\varphi_{\calF}: \pi_{a,*}\calF\otimes\calL^{-1}\to\pi_{a,*}\calF$ because $\calF$ is an $\calO_{Y_{a}}$-module and $\calO_{Y_{a}}$ contains $\calL^{-1}$ as the second direct summand. When $\calF$ is torsion-free and generically rank $1$, $\calV:=\pi_{a,*}\calF$ is torsion-free over $X$ (hence a vector bundle) of rank $n$, and the map $\varphi_{\calF}$ induces a Higgs field $\varphi:\calV\to\calV\otimes\calL$. The assignment $\calF\mapsto (\calV,\varphi)$ defines the morphism $\cPic(Y_{a})\to \calM_{a}$, which can be shown to be an isomorphism. \qed




\subsubsection{The case $G=\Sp_{2n}$} In this case $a\in\calA_{G,\calL}$ is a tuple $(a_{1},\cdots, a_{n})$ with $a_{i}\in\cohog{0}{X,\calL^{\otimes2i}}$. For $a\in\calA_{G,\calL}$, we can similarly define a spectral curve $Y_{a}$ as the closed subscheme of the total space of $\calL$ cut out by the ideal locally generated by
\begin{equation*}
P_{a}(y)=y^{2n}+a_{1}y^{2n-2}+\cdots+a_{n}.
\end{equation*}
Note that $Y_{a}$ carries an involution  $\sigma(y)=-y$ under which the projection $\pi_{a}: Y_{a}\to X$ is invariant. Now suppose $P_{a}(y)$ is separable when restricted to the generic point of $X$, so that $Y_{a}$ is reduced. The involution $\sigma$ induces an involution $\sigma$ on the compactified Picard $\cPic(Y_{a})$. Let $(-)^{\vee}$ be the relative Serre duality functor on coherent sheaves on $Y_{a}$, i.e., $\calF^{\vee}=\un{\Hom}_{\calO_{X}}(\calF,\calO_{X})$ viewed as an $\calO_{Y_{a}}$-module in a natural way. Let $\cPrym(Y_{a};\sigma)$ be the moduli stack of pairs $(\calF,\iota)$ where $\calF\in\cPic(Y_{a})$ and $\iota$ is an isomorphism $\iota: \sigma^{*}\calF\cong\calF^{\vee}$ satisfying that $(\sigma^{*}\iota)^{\vee}=-\iota$. We called $\cPrym(Y_{a};\sigma)$ {\em the compactified Prym stack} of $Y_{a}$ with respect to the involution $\sigma$. Similar to the case of $\GL_{n}$, we have the following description of $\calM_{a}$.

\begin{prop} Suppose $a\in\calA_{G,\calL}$ such that $Y_{a}$ is reduced. Then there is a canonical isomorphism of stacks
\begin{equation*}
\cPrym(Y_{a})\cong\calM_{a}.
\end{equation*}
\end{prop}

\subsubsection{Sketch of proof}
For $a\in\calA_{G,\calL}$, points in $\calM_{a}$ are triples $(\calV,\omega, \varphi)$ where $\calV$ is a vector bundle of rank $2n$ on $X$, $\omega: \wedge^{2}\calV\to\calO_{X}$ is a symplectic form on $\calV$ as in \S\ref{sss:Bun Sp}, and $\varphi:\calV\to\calV\otimes\calL$ satisfies $\om(\varphi u, v)+\om(u, \varphi v)=0$ for local sections $u,v$ of $\calV$, and the characteristic polynomial of $\varphi$ is $P_{a}(y)$. By Proposition \ref{p:cPic}, the Higgs bundle $(\calV,\varphi)$ gives a point $\calF\in\cPic(Y_{a})$. The symplectic form can be viewed as an isomorphism $\jmath: (\calV,\varphi)\isom (\calV^{\vee}, -\varphi^{\vee})$ such that $\jmath^{\vee}=-\jmath$. Note that the Higgs bundle $(\calV^{\vee}, -\varphi^{\vee})$ corresponds to $\sigma^{*}\calF^{\vee}$ under the isomorphism in Proposition \ref{p:cPic}. The isomorphism $\jmath$ then turns into $\iota:  \sigma^{*}\calF\isom\calF^{\vee}$ satisfying $(\sigma^{*}\iota)^{\vee}=-\iota$.  \qed







\subsubsection{Example \ref{sss:MHilb} continued}\label{sss:Hilb}
Fix two line bundles $\calL_{1}$ and $\calL_{2}$ on $X$, and let $\calH_{\calL_{1},\calL_{2}}$ be the fiber of the moduli stack $\calM_{H,\rho}$ in \S\ref{sss:MHilb} over $(\calL_{1},\calL_{2})$. Then $\calH_{\calL_{1},\calL_{2}}$ classifies $(\calV,\varphi,\beta)$ where $\calV$ is a vector bundle of rank $n$ over $X$, $\varphi:\calU\to\calU\otimes\calL_{1}$ is a Higgs field and $\beta: \calU\to \calL_{2}$. We have a Hitchin-type map $\calH_{\calL_{1},\calL_{2}}\to\calA_{n,\calL_{1}}$ sending $(\calV,\varphi,\beta)$ to the $(a_{i}(\varphi))_{1\leq i\leq n}$. Let $a\in\Ah_{n,\calL_{1}}$ and let $\calH_{\calL_{1},\calL_{2},a}$ be the fiber of $\calH_{\calL_{1},\calL_{2}}$ over $a$. Consider the same spectral curve $Y_{a}$ as in \S\ref{sss:spectral curve}. Using Proposition \ref{p:cPic} we may identify $(\calU,\varphi)$ with a point $\calF\in\cPic(Y_{a})$. The map $\beta:\pi_{a,*}\calF=\calU\to\calL_{2}$ gives a map $b:\calF\to \pi^{!}_{a}\calL_{2}$ by adjunction. Here $\pi^{!}_{a}\calL_{2}\cong\pi^{*}_{a}\calL_{2}\otimes\omega_{Y_{a}/X}$, and the relative dualizing complex $\omega_{Y_{a}/X}=\omega_{Y_{a}}\otimes\pi_{a}^{*}\omega_{X}^{-1}$ is a line bundle on $Y_{a}$ because $Y_{a}$ is a planar curve hence Gorenstein. Since $\calF$ is torsion-free and generically a line bundle, $b$ is an injective map of coherent sheaves. Hence the data $(\calU,\varphi,\beta)$ turns into the data of a coherent subsheaf $\calF$ of the line bundle $\pi^{!}_{a}\calL_{2}$. Since $\pi^{!}_{a}\calL_{2}$ is a line bundle, the subsheaf $\calF$ is determined by the support of the quotient $(\pi^{!}_{a}\calL_{2})/\calF$, which is a zero-dimensional subscheme of $Y_{a}$.  We conclude that there is a canonical isomorphism
\begin{equation*}
\calH_{\calL_{1},\calL_{2},a}\cong\Hilb(Y_{a})
\end{equation*}
where $\Hilb(Y_{a})$ is the disjoint union of Hilbert schemes of zero-dimensional subschemes of $Y_{a}$ of various lengths.

To conclude, the Hilbert scheme of points of a family of planar curves can be realized as a Hitchin-type moduli space in the framework of \S\ref{sss:general Hitchin}.




\subsubsection{Symmetry on Hitchin fibers}\label{sss:sym Hit}
In the case $G=\GL_{n}$ and $a\in\Ah_{n,\calL}$, we have identified the Hitchin fiber $\calM_{a}$ with the compactified Picard stack $\cPic(Y_{a})$ of the spectral curve $Y_{a}$. The usual Picard stack $\Pic(Y_{a})$ acts on $\cPic(Y_{a})$ by tensor product. Therefore we have an action of $\Pic(Y_{a})$ on $\calM_{a}$. This action is simply transitive if $Y_{a}$ is smooth.

In the case $G=\Sp_{2n}$ and $a\in\Ah_{G,\calL}$, the Prym stack $\Prym(Y_{a};\sigma)$ is defined as the moduli stack of pairs $(\calL,\iota)$ where $\calL\in\Pic(Y_{a})$ and $\iota:\sigma^*\calL\cong\calL^{-1}$ satisfying $(\sigma^{*}\iota)^{\vee}=\iota$. Then $\Prym(Y_{a};\sigma)$ acts on $\calM_{a}\cong\cPrym(Y_{a};\sigma)$ by tensoring.

For general $G$ and $a\in\Ah_{G,\calL}$, one can similarly define a commutative group stack $\calP_{a}$ that acts on the corresponding Hitchin fiber $\calM_{a}$. In fact, via the map $a:X\to [\frc/\Gm]$, the regular centralizer group scheme $J$ in \S\ref{sss:sym ASF}, which descends to $[\frc/\Gm]$, pulls back to a smooth group scheme $J_{a}$ over $X$ which is generically a torus. The stack $\calP_{a}$ then classifies $J_{a}$-torsors over $X$. 


\subsection{Relation with affine Springer fibers}\label{ss:rel Hit ASF}
In this subsection we will state a precise relationship between Hitchin fibers and affine Springer fibers, observed by Ng\^o. 

\subsubsection{} Fix a point $a\in\Ah_{G,\calL}(k)$. Viewing $a$ as a map $X\to [\frc/\Gm]$, let $U_{a}\subset X$ be the preimage of $[\frc^{\rs}/\Gm]$ under $a$. Since $a\in\Ah_{G,\calL}$, $U_{a}$ is non-empty hence the complement $X-U_{a}$ consists of finitely many $k$-points of $X$.  For each point $x\in X-U_{a}$, let $\calO_{x}$ be the completed local ring of $X$ at $x$, with fraction field $F_{x}$ and residue field $k(x)=k$.  We fix a trivialization of $\calL$ near $x$ and we may identify $a_{x}=a|_{\Spec\calO_{x}}$ as an element in $\frc(\calO_{x})\cap\frc^{\rs}(F_{x})$. Let $\ep(a_{x})\in\frg^{\rs}(F_{x})$ be the corresponding point in the Kostant section, then we have the affine Springer fiber $\sX_{a_{x}}:=\sX_{\ep(a_{x})}$ in the affine Grassmannian $\Gr_{G, x}=L_{x}G/L^{+}_{x}G$ (we put $x$ in the subscript to emphasize that the definition of the loop groups uses the field $F_{x}$, which is isomorphic to $k\lr{t}$). The loop group of the centralizer $LG_{\ep(a_{x})}$ acts on $\sX_{a_{x}}$, and the action  factors through the local Picard group $P_{a_{x}}$ as in \S\ref{sss:sym ASF}. On the other hand, we have the action of the global Picard stack $\calP_{a}$ on the Hitchin fiber $\calM_{a}$ mentioned in \S\ref{sss:sym Hit}. The product formula of Ng\^o roughly says that, modulo the actions of the local and global Picard stacks, $\calM_{a}$ is the product of the affine Springer fibers $\sX_{a_{x}}$ for all points $x\in X-U_{a}$.

\begin{theorem}[Product Formula, Ng\^o {\cite[Th 4.6]{NgoHit} and \cite[Prop 4.15.1]{NgoFL}}]\label{th:prod} For $a\in\Ah(k)$, there is a canonical morphism
\begin{equation*}
\calP_{a}\twtimes{\prod_{x\in X-U_{a}}P_{a_{x}}}\left(\prod_{x\in X-U_{a}}\sX_{a_{x}}\right)\to \calM_{a}
\end{equation*}
which is a homeomorphism of stacks. Here the notation $P\twtimes{H}Y$ (where $H$ acts on $P$ on the right and acts on $Y$ on the left) means the quotient of $P\times Y$ by the action of $H$ given by $h\cdot (p,y)=(ph^{-1}, hy)$, $p\in P, y\in Y$ and $h\in H$.
\end{theorem}

In the case $G=\GL_{n}$, the product formula can be reinterpreted in more familiar terms using the compactified Picard stack of spectral curves, which in fact makes sense for all reduced curves. Let $C$ be a reduced and projective curve over $k$. For each point $x\in C(k)$, one can define a local Picard group $P_{x}(C)$ whose $k$ points are $F^{\times}_{x}/\calO^{\times}_{x}$, where $\calO_{x}$ is the completed local ring of $C$ at $x$ and $F_{x}$ its ring of total fractions (which is a product of fields in general). There is also a local analog $\ov{P}_{x}(C)$ of $\cPic(C)$ whose $k$-points are the fractional $\calO_{x}$-ideals of $F_{x}$, compare \S\ref{sss:GLn orb formula}. Then the following variant of the product formula holds, whose proof is similar to that of Theorem \ref{th:prod}.

\begin{prop}\label{p:Pic prod} Let $C$ be a reduced and projective curve over $k$. Let $Z=C-C^{sm}$ be the singular locus of $C$. Then there is a canonical morphism 
\begin{equation*}
\Pic(C)\twtimes{\prod_{x\in Z}P_{x}(C)}\left(\prod_{x\in Z}\ov{P}_{x}(C)\right)\to\cPic(C)
\end{equation*}
which is a homeomorphism of stacks.
\end{prop}

The product formula provides a link between the geometry of Hitchin fibers and that of affine Springer fibers, the latter is closely related to orbital integrals as we have already seen. This link makes it possible to approach the Fundamental Lemma by studying the cohomology of Hitchin fibers. The advantage of using Hitchin fibers instead of affine Springer fibers is that  the Hitchin fibration has nicer geometric properties, as we have seen in \S\ref{sss:geom Hit}.

